% 02-apprentissage-mecanique.tex — État de l'art 1/5 : le cours de Physique 1,
% les conceptions erronées en dynamique, la mesure de la compréhension et ce
% qu'elles demandent à une animation.

\chapter{Apprendre la mécanique en première année}\label{chap:mecanique}

Ce chapitre résume ce que la recherche en didactique de la physique sait des
difficultés des étudiants en mécanique. Il en tire les premières exigences
pour nos animations.

\section{Le cours de Physique~1 (BAC1 ingénieur civil)}

Le cours LEPL1201 couvre la mécanique du point et du solide, puis
l'électricité élémentaire. Son formulaire, autorisé à l'examen, fixe le
contenu et les notations \cite{formulaire_epl1201} :
\begin{itemize}
  \item la cinématique, dont le mouvement d'un projectile écrit composante
        par composante : \(x(t) = u_0 t + x_0\) et
        \(y(t) = -g t^2/2 + v_0 t + y_0\), avec le vocabulaire du mouvement
        rectiligne uniforme (MRU) et uniformément accéléré (MRUA) ;
  \item le mouvement circulaire, avec \(v = r \omega\), l'accélération
        radiale \(a_r = -r \omega^2\) et tangentielle \(a_\theta = r \alpha\) ;
  \item la dynamique (lois de Newton, moment cinétique) et le théorème de
        l'énergie mécanique ;
  \item la force de gravitation et la force électrique, le théorème de Gauss
        et le potentiel électrique.
\end{itemize}
Ces notations comptent pour une animation. Un étudiant qui voit \(v_{0x}\)
dans une simulation et \(u_0\) au cours doit faire une traduction de plus,
qui n'apporte rien et peut l'induire en erreur. Toutes nos pages reprennent
donc les symboles du formulaire.

\section{Conceptions erronées en dynamique}

Les étudiants n'arrivent pas en mécanique sans idées. Ils ont une
\enquote{physique du sens commun}, construite par l'expérience quotidienne,
et elle contredit souvent les lois de Newton. Halloun et Hestenes l'ont
mesurée chez des étudiants d'université avant et après un cours classique :
beaucoup pensent qu'un mouvement exige une force dans le sens de la vitesse,
et cette idée change peu après le cours \cite{halloun1985common}. Viennot
avait observé le même raisonnement chez des étudiants de plusieurs pays
\cite{viennot1979}. On retrouve dans ces travaux, et dans l'expérience
courante des enseignants, les conceptions suivantes :
\begin{itemize}
  \item une force est nécessaire pour entretenir le mouvement, et elle pointe
        dans le sens de la vitesse ;
  \item la vitesse et l'accélération ont la même direction, en particulier
        dans un virage ;
  \item dans un mouvement circulaire, une \enquote{force centrifuge} pousse
        l'objet vers l'extérieur ;
  \item la réaction normale est toujours égale au poids.
\end{itemize}

Le mouvement circulaire réunit plusieurs de ces difficultés. L'accélération
pointe vers le centre alors que la vitesse est tangente, ce qui contredit
l'idée que l'accélération suit la vitesse. Il n'existe pas de \enquote{force
centripète} qui s'ajouterait aux autres : c'est la résultante du poids, de la
réaction normale ou de la tension qui joue ce rôle.

Le looping en est l'exemple classique. Prenons un wagon assimilé à un point,
lâché sans vitesse à la hauteur \(h\) au-dessus du bas d'un looping
circulaire de rayon \(R\), sans frottement. Pour savoir s'il fait le tour, il
faut combiner deux principes. La conservation de l'énergie donne sa vitesse
au sommet, \(v_s^2 = 2 g (h - 2R)\). La deuxième loi de Newton selon le rayon
donne alors la réaction du rail au sommet, \(N_s = m (v_s^2/R - g)\), qui doit
rester positive ou nulle : sinon, le rail devrait tirer le wagon, ce qu'il ne
peut pas faire. On obtient la hauteur minimale \(h_{\min} = 5R/2\). Le seau
d'eau que l'on fait tourner dans un plan vertical obéit à la même condition,
avec la tension de la corde à la place de la réaction normale. Les vrais
loopings ne sont d'ailleurs pas circulaires, pour limiter l'accélération
subie par les passagers \cite{pendrill2005loops}, ce qui ouvre l'exercice sur
une situation réelle.

\section{Mesurer la compréhension conceptuelle}

Pour mesurer ces conceptions, Hestenes, Wells et Swackhamer ont construit le
\emph{Force Concept Inventory} (FCI) \cite{hestenes1992fci}. C'est un
questionnaire à choix multiples sans calcul, dont chaque mauvaise réponse
correspond à une conception erronée précise. Il est devenu l'instrument de
référence pour comparer des méthodes d'enseignement de la mécanique.

Hake l'a utilisé sur plus de six mille étudiants, dans des cours d'écoles
secondaires, de collèges et d'universités \cite{hake1998}. Pour comparer des
groupes qui partent de niveaux différents, il utilise le \emph{gain
normalisé} : la part de ce qui restait à apprendre qui a été acquise pendant
le cours,
\[
  \langle g \rangle = \frac{\text{score final} - \text{score initial}}
                           {100\,\% - \text{score initial}} .
\]
Ce gain vaut en moyenne \num{0,23} dans les cours magistraux classiques et
\num{0,48} dans les cours à \enquote{engagement interactif}, où l'étudiant
est actif et reçoit un retour immédiat sur ses raisonnements. Le constat est
net : un cours magistral, même bien donné, corrige peu les conceptions
erronées ; il faut que l'étudiant confronte ses idées à ce qui se passe
réellement.

Ce type de mesure servira à évaluer nos animations si un test avec des
étudiants est possible (chapitre~\ref{chap:protocole}).

\section{Ce que cela demande à une animation}

Ces travaux donnent trois exigences directes :
\begin{itemize}
  \item montrer \emph{toutes} les forces, avec leur direction et leur
        grandeur, pour que l'étudiant voie qu'aucune force centrifuge
        n'apparaît, et que la réaction normale varie ;
  \item montrer la vitesse et l'accélération séparément, pour que leur
        différence de direction soit visible ;
  \item laisser l'étudiant prédire puis vérifier : changer un paramètre, voir
        si le wagon fait le tour, et comprendre pourquoi.
\end{itemize}
À ces exigences s'ajoute celle de la section précédente : parler la langue du
cours, avec ses notations.
